% ============================================================================
% COURS 12 — EXERCICES, PARTIE B
% Questions de cours, autocontraintes et QCM
% Source : 16_Rdm_2026.docx
% ============================================================================

\subsection{Questions de cours}

\begin{enumerate}
  \item Préciser la signification de l'acronyme RdM et citer les deux objectifs principaux d'une étude de résistance des matériaux.
  \item Citer les cinq principales hypothèses qui permettent de simplifier une étude de poutre.
  \item Écrire le torseur de cohésion complet et donner le nom de ses six composantes.
  \item Donner la définition du vecteur contrainte $\vec t(M,\vec n)$.
  \item Écrire la condition de résistance d'une pièce soumise à la traction et préciser tous les termes utilisés.
  \item Citer six sollicitations simples ou usuelles.
  \item Écrire les lois de Hooke en contrainte normale et en cisaillement. Donner le nom des modules $E$ et $G$ ainsi que leur relation pour un matériau isotrope :
  \[
  G=\frac{E}{2(1+\nu)}.
  \]
  \item Compléter le vocabulaire associé aux symboles suivants : $\nu$, $\theta$, $\varepsilon$, $\rho$, $\sigma$, $\tau$ et $\gamma$.
  \item Associer chaque sollicitation à la déformation et au torseur de cohésion correspondants.
  \item Dans un schéma de chargement, préciser le nom de l'action $F$ et son unité.
  \item Écrire l'expression de la contrainte normale en traction puis en flexion.
  \item Donner le nom et l'unité de $I_{Gz}$, $y$, $M_{fz}$ et $R_{pe}$.
  \item Donner l'équation différentielle de la déformée en flexion.
  \item Donner les conditions aux limites usuelles :
  \begin{itemize}
    \item encastrement en $x=0$ ;
    \item appui simple en $x=a$ ;
    \item continuité entre deux tronçons au point $x=a$.
  \end{itemize}
\end{enumerate}

\subsection{Défauts géométriques et autocontraintes}

L'arbre 2 est lié au carter 1 par deux roulements. En $O_1$, un roulement à deux rangées de billes à contacts obliques est modélisé par une liaison pivot d'axe $\vec x_1$. En $A$, un roulement à billes à contact radial dont la bague extérieure est libre axialement est modélisé par une liaison linéaire annulaire.

Après fabrication, les deux alésages du carter présentent un défaut de coaxialité radial $\delta_r$. Le montage force l'arbre à se déformer et crée des actions de liaison sans qu'aucune charge extérieure ne soit appliquée : ce sont des \textbf{autocontraintes}.

\begin{center}
\IfFileExists{pandoc/media/image80.png}{\includegraphics[width=.70\linewidth]{pandoc/media/image80.png}}{}
\end{center}

On suppose le problème plan. Le plan est choisi de sorte que la composante hors plan du défaut soit nulle. Les composantes des torseurs de liaison sont notées $X_O$, $Y_O$, $N_O$ au point $O_1$ et $Y_A$ au point $A$.

\subsubsection{Étude statique et hyperstatisme}

\begin{enumerate}
  \item Écrire, dans le problème plan, les torseurs des actions de liaison du carter sur l'arbre aux points $O_1$ et $A$.
  \item Faire le bilan des inconnues et déterminer le degré d'hyperstatisme interne de la liaison équivalente entre l'arbre et le carter.
  \item Aucune charge extérieure n'agissant sur l'arbre, appliquer le PFS et établir les trois relations entre les inconnues de liaison.
\end{enumerate}

\subsubsection{Déformée de l'arbre}

La section de l'arbre possède un moment quadratique $I_{Gz}$ et le matériau un module d'Young $E$.

\begin{enumerate}[resume]
  \item Montrer que le moment fléchissant, pour $0\leq x\leq L$, peut s'écrire :
  \[
  M_{fz}(x)=-N_O+xY_O.
  \]
  \item En utilisant les conditions d'encastrement équivalent au point $O_1$, déterminer l'expression de la déformée $y(x)$.
  \item Utiliser la condition géométrique $y(L)=\delta_r$ pour lever l'hyperstaticité et montrer que :
  \[
  \boxed{Y_A=\frac{3EI_{Gz}\delta_r}{L^3}.}
  \]
  Déterminer ensuite $Y_O$ et $N_O$.
  \item En déduire l'expression complète de $y(x)$ en fonction de $E$, $I_{Gz}$, $L$ et $\delta_r$.
\end{enumerate}

\subsubsection{Application numérique}

On donne :
\[
D=20\ \mathrm{mm},\qquad L=300\ \mathrm{mm},\qquad
E=2\times10^5\ \mathrm{MPa},\qquad \delta_r=0{,}05\ \mathrm{mm}.
\]

\begin{enumerate}[resume]
  \item Calculer le moment quadratique :
  \[
  I_{Gz}=\frac{\pi D^4}{64}.
  \]
  \item Calculer les actions de liaison.
  \item Déterminer le moment fléchissant maximal puis l'autocontrainte normale maximale dans l'arbre.
\end{enumerate}

\subsection{QCM de synthèse}

Une seule réponse est attendue pour chaque question. Pour les applications numériques, on pourra prendre $\pi=3$ et $g=10\ \mathrm{m/s^2}$ lorsque cela simplifie les calculs.

\begin{enumerate}
  \item Comment s'appelle le diagramme représentant la contrainte en fonction de la déformation lors d'un essai de traction ?
  \begin{itemize}
    \item loi de Hooke ; essai de traction ; diagramme contrainte--déformation ; module d'Young.
  \end{itemize}

  \item Sur une courbe de traction conventionnelle, comment nomme-t-on la zone comprise entre la limite d'élasticité et la rupture ?
  \begin{itemize}
    \item domaine élastique ; domaine plastique ; striction ; écrouissage.
  \end{itemize}

  \item Parmi les affirmations suivantes, laquelle caractérise correctement un acier dont la limite d'élasticité vaut $295\ \mathrm{MPa}$ ?
  \begin{itemize}
    \item au-delà de cette contrainte, les déformations peuvent devenir permanentes ;
    \item la contrainte est alors nécessairement tangentielle ;
    \item la loi $\sigma=E\varepsilon$ reste valable jusqu'à la rupture.
  \end{itemize}

  \item Quel est l'ordre de grandeur du module d'Young du plexiglas ?
  \begin{itemize}
    \item $3{,}1\ \mathrm{GPa}$ ; $62\ \mathrm{GPa}$ ; $210\ \mathrm{GPa}$.
  \end{itemize}

  \item Déterminer la coordonnée du centre d'aire de la section composée représentée ci-dessous.
  \begin{center}
  \IfFileExists{pandoc/media/image82.png}{\includegraphics[width=.30\linewidth]{pandoc/media/image82.png}}{}
  \end{center}

  \item Quelle affirmation concernant le torseur de cohésion est fausse ?
  \begin{itemize}
    \item $M_f$ désigne un moment fléchissant ;
    \item $T_z$ est l'effort tranchant suivant $\vec z$ ;
    \item $N$ est l'effort tranchant suivant $\vec y$ ;
    \item $M_t$ désigne le moment de torsion.
  \end{itemize}

  \item Une poutre dont le torseur de cohésion ne comporte que $M_{fz}$ est soumise à :
  \begin{itemize}
    \item traction ; compression ; cisaillement ; torsion ; flexion pure ; flexion simple.
  \end{itemize}

  \item Quel terme n'est pas une grandeur ou une notion usuelle de RdM ?
  \begin{itemize}
    \item résistance mécanique ; contrainte tangentielle ; moment fléchissant ; module de Coulomb ; pression uniforme.
  \end{itemize}

  \item Dans un torseur de cohésion, que représente $N$ ?
  \begin{itemize}
    \item un effort normal ; un moment nominal ; une unité en newtons.
  \end{itemize}

  \item Identifier la sollicitation principale de la pièce représentée.
  \begin{center}
  \IfFileExists{pandoc/media/image84.png}{
\includegraphics[width=.48\linewidth]{pandoc/media/image84.png}}{}
  \end{center}

  \item Associer chacune des déformations usuelles au torseur de cohésion réduit correspondant : traction, cisaillement, torsion, flexion pure et flexion simple.

  \item En quelle unité s'exprime un moment quadratique ou un moment polaire de section ?
  \begin{itemize}
    \item $\mathrm{mm^2}$ ; $\mathrm{mm^3}$ ; $\mathrm{mm^4}$ ; $\mathrm{N/mm^2}$.
  \end{itemize}

  \item Une éprouvette reçoit une force normale de $52\ \mathrm{kN}$ sur une section de $260\ \mathrm{mm^2}$. Son module d'Young vaut $40\ \mathrm{GPa}$. Calculer la déformation relative.
  \begin{itemize}
    \item $0{,}3\ \%$ ; $0{,}5\ \%$ ; $1{,}5\ \%$ ; $2\ \%$.
  \end{itemize}

  \item Quelle lettre désigne le premier coefficient de Lamé ?
  \begin{itemize}
    \item $\nu$ ; $\varepsilon$ ; $\lambda$ ; $\theta$ ; $\rho$ ; $\tau$ ; $\gamma$ ; $\sigma$.
  \end{itemize}

  \item Dans une section droite soumise à de la flexion pure, le torseur de cohésion est :
  \begin{itemize}
    \item nul ; un glisseur ; un couple ; un torseur quelconque.
  \end{itemize}

  \item Une force normale de $48\ \mathrm N$ est uniformément répartie sur une section circulaire de rayon $12\ \mathrm{cm}$. Calculer la contrainte moyenne.
  \begin{itemize}
    \item $120\ \mathrm{Pa}$ ; $1{,}1\ \mathrm{kPa}$ ; $6{,}7\ \mathrm{kPa}$ ; $280\ \mathrm{kPa}$.
  \end{itemize}

  \item Une poutre horizontale de longueur $L$, encastrée à une extrémité, reçoit une force verticale $F$ à l'autre extrémité. Quel est le moment fléchissant maximal ?
  \begin{itemize}
    \item $FL$ ; $FL/2$ ; $FL/4$ ; $FL^2/2$.
  \end{itemize}

  \item Une section circulaire de diamètre $100\ \mathrm{mm}$ reçoit un moment fléchissant de $75\ \mathrm{N\,m}$. Calculer la contrainte normale maximale.
  \begin{itemize}
    \item $160\ \mathrm{Pa}$ ; $400\ \mathrm{kPa}$ ; $0{,}8\ \mathrm{MPa}$ ; $12{,}8\ \mathrm{MPa}$.
  \end{itemize}

  \item Comment modélise-t-on le poids propre d'une poutre droite homogène ?
  \begin{itemize}
    \item par un glisseur unique au centre de gravité ;
    \item par une densité linéique de force ;
    \item par des moments appliqués aux appuis.
  \end{itemize}

  \item Parmi les propositions suivantes, laquelle correspond à une sollicitation composée ?
  \begin{itemize}
    \item extension ; cisaillement ; flexion déviée ; torsion simple.
  \end{itemize}

  \item Dans un essai idéal de cisaillement suivant $\vec y$, quelle composante parasite du torseur de cohésion doit être négligée pour parler de cisaillement simple ?
  \begin{itemize}
    \item $N$ ; $T_y$ ; $T_z$ ; $M_t$ ; $M_{fy}$ ; $M_{fz}$.
  \end{itemize}

  \item Une poutre circulaire de rayon $12\ \mathrm{mm}$ reçoit un moment de torsion de $9\ \mathrm{N\,m}$. Calculer la contrainte tangentielle maximale avec :
  \[
  \tau_{\max}=\frac{M_tR}{J},
  \qquad
  J=\frac{\pi R^4}{2}.
  \]
  \begin{itemize}
    \item $3{,}5\ \mathrm{kPa}$ ; $3{,}5\ \mathrm{MPa}$ ; $6{,}9\ \mathrm{MPa}$ ; $27{,}8\ \mathrm{MPa}$.
  \end{itemize}
\end{enumerate}
